\section{September 17, 2026}

We continue the study of measurable functions and then introduce the
building blocks of the Lebesgue integral.  Throughout, $(X,\mathcal M)$
is a measurable space.  A real- or complex-valued function is understood
to have the Borel $\sigma$-algebra on its codomain.

Recall that the Borel $\sigma$-algebra on
$\overline{\bb R}=\bb R\cup\{-\infty,+\infty\}$ is generated by either
of the families
\[
  \{(a,+\infty]:a\in\bb R\}
  \quad\text{and}\quad
  \{[-\infty,a]:a\in\bb R\}.
\]
Thus the half-line criterion from the preceding lecture applies equally
well to extended-real-valued functions.

\subsection{Operations on measurable functions}

Let $(Y_i,\mathcal N_i)_{i\in I}$ be measurable spaces, and give
\[
  Y=\prod_{i\in I}Y_i
\]
the product $\sigma$-algebra
$\mathcal N=\bigotimes_{i\in I}\mathcal N_i$.  If
$\pi_i\colon Y\to Y_i$ is the $i$th coordinate projection, then
$\mathcal N$ is generated by the cylinder sets
\[
  \pi_i^{-1}(E),
  \qquad i\in I,\quad E\in\mathcal N_i.
\]

\begin{lemma}[Coordinate criterion for product measurability]
  \label{lem:product-coordinate-measurability-sep-17}
  A map $f\colon X\to Y$ is
  $(\mathcal M,\bigotimes_{i\in I}\mathcal N_i)$-measurable if and only if
  every coordinate map
  \[
    \pi_i\circ f\colon X\longrightarrow Y_i
  \]
  is $(\mathcal M,\mathcal N_i)$-measurable.
\end{lemma}

\begin{proof}
  If $f$ is measurable, then $\pi_i\circ f$ is measurable because each
  projection $\pi_i$ is measurable.  Conversely, suppose all the
  coordinate maps are measurable.  For every generating cylinder,
  \[
    f^{-1}\bigl(\pi_i^{-1}(E)\bigr)
    =(\pi_i\circ f)^{-1}(E)\in\mathcal M.
  \]
  Lemma~\ref{lem:measurability-generators-sep-15} now shows that $f$ is
  measurable.
\end{proof}

\begin{figure}[htbp]
  \centering
  \begin{tikzpicture}[
    node distance=2.1cm and 3.0cm,
    every node/.style={font=\small},
    map/.style={-{Latex[length=2.2mm]},thick}]
    \node (X) {$(X,\mathcal M)$};
    \node[right=of X] (Y) {$\left(\prod_{i\in I}Y_i,
      \bigotimes_{i\in I}\mathcal N_i\right)$};
    \node[below=of Y] (Yi) {$(Y_i,\mathcal N_i)$};
    \draw[map] (X)--node[above] {$f$}(Y);
    \draw[map] (Y)--node[right] {$\pi_i$}(Yi);
    \draw[map] (X)--node[below left] {$f_i=\pi_i\circ f$}(Yi);
  \end{tikzpicture}
  \caption{Measurability into a product can be checked one coordinate
    at a time.}
  \label{fig:product-measurability-sep-17}
\end{figure}

\begin{corollary}\label{cor:complex-coordinate-measurability-sep-17}
  A function $f\colon X\to\bb C$ is measurable if and only if
  $\operatorname{Re}f$ and $\operatorname{Im}f$ are measurable.
\end{corollary}

\begin{proof}
  Under the identification $\bb C\cong\bb R^2$,
  \[
    \mathcal B_{\bb C}
    =\mathcal B_{\bb R}\otimes\mathcal B_{\bb R}.
  \]
  The claim is therefore the two-coordinate case of
  Lemma~\ref{lem:product-coordinate-measurability-sep-17}.
\end{proof}

\begin{proposition}[Algebra and lattice operations]
  \label{prop:measurable-function-operations-sep-17}
  If $f,g\colon X\to\bb R$ are measurable, then so are
  \[
    f+g,\quad fg,\quad f^+=\max\{f,0\},\quad
    f^-=\max\{-f,0\},\quad |f|,
  \]
  and $\operatorname{sgn}(f)$.
\end{proposition}

\begin{proof}
  By Lemma~\ref{lem:product-coordinate-measurability-sep-17}, the map
  \[
    F=(f,g)\colon X\longrightarrow\bb R^2
  \]
  is measurable.  Addition and multiplication are continuous maps from
  $\bb R^2$ to $\bb R$, so they are Borel measurable.  Hence
  \[
    f+g=\Sigma\circ F,
    \qquad
    fg=\Pi\circ F,
  \]
  where $\Sigma(s,t)=s+t$ and $\Pi(s,t)=st$, are measurable.  The
  remaining functions follow by composition with the Borel maps
  \[
    t\longmapsto\max\{t,0\},\quad
    t\longmapsto\max\{-t,0\},\quad
    t\longmapsto |t|,
  \]
  and with the Borel measurable sign function.  Notice also that
  $f=f^+-f^-$ and $|f|=f^++f^-$.
\end{proof}

The same argument, together with
Corollary~\ref{cor:complex-coordinate-measurability-sep-17}, shows that
sums and products of complex-valued measurable functions are measurable.

\begin{proposition}[Countable suprema and limits]
  \label{prop:measurable-sequential-operations-sep-17}
  Let $f_n\colon X\to\overline{\bb R}$ be measurable for every
  $n\in\bb N$.  Then the functions
  \[
    \sup_n f_n,\qquad \inf_n f_n,\qquad
    \limsup_{n\to\infty}f_n,
    \qquad \liminf_{n\to\infty}f_n
  \]
  are measurable.  In particular, if $f_n\to f$ pointwise, then $f$ is
  measurable.
\end{proposition}

\begin{proof}
  For $a\in\bb R$,
  \[
    \left\{\sup_n f_n>a\right\}
      =\bigcup_{n=1}^{\infty}\{f_n>a\},
    \qquad
    \left\{\inf_n f_n<a\right\}
      =\bigcup_{n=1}^{\infty}\{f_n<a\}.
  \]
  The half-line criterion implies that the pointwise supremum and
  infimum are measurable.  Now use the identities
  \[
    \limsup_{n\to\infty}f_n
      =\inf_{N\geq1}\sup_{n\geq N}f_n,
    \qquad
    \liminf_{n\to\infty}f_n
      =\sup_{N\geq1}\inf_{n\geq N}f_n.
  \]
  If $f_n\to f$ pointwise, both expressions equal $f$.
\end{proof}

\begin{figure}[htbp]
  \centering
  \begin{tikzpicture}[x=1.02cm,y=1.05cm,>=Latex]
    \draw[->] (0.55,0.45)--(9.25,0.45) node[right] {$n$};
    \draw[->] (0.75,0.25)--(0.75,4.0) node[above] {value at a fixed $x$};
    \draw[gray!55,dashed] (0.75,2.2)--(9,2.2)
      node[right,black] {$\lim f_n(x)$};

    \foreach \n/\v in {1/3.45,2/1.0,3/3.05,4/1.4,5/2.72,6/1.72,
      7/2.48,8/2.2} {
      \fill[black] (\n,\v) circle (1.6pt);
    }
    \draw[blue!75!black,very thick]
      plot coordinates {(1,3.45)(2,3.05)(3,3.05)(4,2.72)
        (5,2.72)(6,2.48)(7,2.48)(8,2.2)};
    \draw[red!75!black,very thick]
      plot coordinates {(1,1.0)(2,1.0)(3,1.4)(4,1.4)
        (5,1.72)(6,1.72)(7,1.95)(8,2.2)};
    \node[blue!75!black,above] at (2.25,3.05)
      {$U_N=\sup_{n\geq N}f_n(x)$};
    \node[red!75!black,below] at (2.25,1.0)
      {$L_N=\inf_{n\geq N}f_n(x)$};
  \end{tikzpicture}
  \caption{The tail supremum decreases and the tail infimum increases.
    Their limits are the pointwise $\limsup$ and $\liminf$; the drawing
    shows the convergent case.}
  \label{fig:limsup-liminf-sep-17}
\end{figure}

\subsection{Simple functions}

For $E\in\mathcal M$, continue to write $\mathbf 1_E$ for the
characteristic function of $E$.

\begin{definition}
  A measurable function $\phi\colon X\to\bb C$ is \emph{simple} if it
  has finite range.  Equivalently, it can be written as a finite linear
  combination
  \[
    \phi=\sum_{j=1}^{N}a_j\mathbf 1_{E_j},
    \qquad a_j\in\bb C,\quad E_j\in\mathcal M.
  \]
  If $\{a_1,\ldots,a_N\}$ is the set of distinct values of $\phi$ and
  $E_j=\phi^{-1}(\{a_j\})$, then the sets $E_j$ are pairwise disjoint
  and cover $X$.  This is the \emph{standard representation} of
  $\phi$.
\end{definition}

Sums, products, absolute values, maxima, and minima of simple functions
are again simple.  The standard representation makes a simple function
look like a finite collection of horizontal levels over a measurable
partition of its domain.

\begin{figure}[htbp]
  \centering
  \begin{tikzpicture}[x=1.35cm,y=0.9cm,>=Latex]
    \draw[->] (-0.1,0)--(7.1,0) node[right] {$X=[0,7]$};
    \draw[->] (0,-0.25)--(0,4.65) node[above] {$\phi$};
    \fill[blue!12] (0,0) rectangle (1.4,1.2);
    \fill[green!16] (1.4,0) rectangle (3.1,3.0);
    \fill[orange!20] (3.1,0) rectangle (4.4,1.9);
    \fill[purple!13] (4.4,0) rectangle (7,3.8);
    \draw[blue!75!black,very thick]
      (0,1.2)--(1.4,1.2) (1.4,3.0)--(3.1,3.0)
      (3.1,1.9)--(4.4,1.9) (4.4,3.8)--(7,3.8);
    \foreach \x in {1.4,3.1,4.4} {
      \draw[gray!55,dashed] (\x,0)--(\x,4.05);
    }
    \node[below] at (0.7,0) {$E_1$};
    \node[below] at (2.25,0) {$E_2$};
    \node[below] at (3.75,0) {$E_3$};
    \node[below] at (5.7,0) {$E_4$};
    \node[left] at (0,1.2) {$a_1$};
    \node[left] at (0,3.0) {$a_2$};
    \node[left] at (0,1.9) {$a_3$};
    \node[left] at (0,3.8) {$a_4$};
  \end{tikzpicture}
  \caption{A simple function
    $\phi=\sum_{j=1}^{4}a_j\mathbf 1_{E_j}$ in standard form.  The
    sets $E_j$ need not be intervals in a general measurable space.}
  \label{fig:simple-function-sep-17}
\end{figure}

Simple functions are sufficient to approximate every measurable
function.

\begin{theorem}[Approximation by simple functions]
  \label{thm:simple-approximation-sep-17}
  \begin{enumerate}[label=(\roman*)]
    \item If $f\colon X\to[0,+\infty]$ is measurable, then there are
      nonnegative simple functions $\phi_n$ such that
      \[
        0\leq\phi_1\leq\phi_2\leq\cdots\leq f,
        \qquad
        \phi_n(x)\longrightarrow f(x)
      \]
      for every $x\in X$.  The convergence is uniform on every set on
      which $f$ is bounded.
    \item If $f\colon X\to\bb C$ is measurable, then there are simple
      functions $\phi_n\colon X\to\bb C$ such that
      \[
        0\leq|\phi_1|\leq|\phi_2|\leq\cdots\leq|f|,
        \qquad
        \phi_n(x)\longrightarrow f(x)
      \]
      for every $x\in X$.  Again, convergence is uniform on every set
      on which $|f|$ is bounded.
  \end{enumerate}
\end{theorem}

\begin{proof}
  For (i), divide the interval $[0,n]$ into dyadic pieces of length
  $2^{-n}$.  Set
  \[
    E_{n,k}
      =f^{-1}\bigl((k2^{-n},(k+1)2^{-n}]\bigr),
      \qquad 0\leq k<n2^n,
  \]
  and $F_n=f^{-1}((n,+\infty])$.  Define
  \begin{equation}\label{eq:dyadic-simple-approximation-sep-17}
    \phi_n
      =\sum_{k=0}^{n2^n-1}k2^{-n}\mathbf 1_{E_{n,k}}
       +n\mathbf 1_{F_n}.
  \end{equation}
  Each $\phi_n$ is a nonnegative simple function.  The dyadic grids
  refine one another and the truncation height increases, so
  $\phi_n\leq\phi_{n+1}\leq f$.  Moreover,
  \[
    0\leq f(x)-\phi_n(x)\leq2^{-n}
    \quad\text{whenever }f(x)\leq n.
  \]
  It follows that $\phi_n\uparrow f$ pointwise.  If $f\leq M$ on a set
  $E$, then for $n\geq M$,
  \[
    \sup_{x\in E}|f(x)-\phi_n(x)|\leq2^{-n},
  \]
  which proves the asserted uniform convergence.

  For (ii), write
  \[
    f=(\operatorname{Re}f)^+-(\operatorname{Re}f)^-
      +i(\operatorname{Im}f)^+-i(\operatorname{Im}f)^-.
  \]
  Apply (i) to these four nonnegative measurable functions to obtain
  increasing simple approximants $u_n,v_n,w_n,z_n$, respectively, and
  set
  \[
    \phi_n=(u_n-v_n)+i(w_n-z_n).
  \]
  At each point, at most one of $u_n,v_n$ is nonzero and at most one of
  $w_n,z_n$ is nonzero.  Consequently $|\phi_n|$ increases with $n$ and
  is bounded by $|f|$.  The convergence statements follow from those
  for the four real components.
\end{proof}

\begin{figure}[!htbp]
  \centering
  \begin{tikzpicture}[x=1.45cm,y=1.35cm,>=Latex]
    \draw[->] (-0.1,0)--(5.15,0) node[right] {$x$};
    \draw[->] (0,-0.1)--(0,3.65) node[above] {$f(x)$};
    \draw[gray!30,dashed] (0,0.5)--(4.8,0.5);
    \draw[gray!30,dashed] (0,1.0)--(4.8,1.0);
    \draw[gray!30,dashed] (0,1.5)--(4.8,1.5);
    \draw[gray!30,dashed] (0,2.0)--(4.8,2.0);
    \draw[gray!30,dashed] (0,2.5)--(4.8,2.5);
    \draw[gray!30,dashed] (0,3.0)--(4.8,3.0);

    \draw[black,very thick,domain=0.2:4.6,samples=80]
      plot (\x,{0.48*(\x-2.35)^2+0.62});
    \draw[blue!75!black,very thick]
      (0.2,2.5)--(0.55,2.5)--(0.55,2.0)--(0.9,2.0)--
      (0.9,1.5)--(1.35,1.5)--(1.35,1.0)--(1.9,1.0)--
      (1.9,0.5)--(2.9,0.5)--(2.9,1.0)--(3.45,1.0)--
      (3.45,1.5)--(3.85,1.5)--(3.85,2.0)--(4.2,2.0)--
      (4.2,2.5)--(4.6,2.5);
    \draw[red!70!black,<->] (3.62,1.5)--(3.62,1.91)
      node[right,font=\small] {$\leq2^{-n}$};
    \node[black] at (4.42,3.3) {$f$};
    \node[blue!75!black] at (4.55,2.25) {$\phi_n$};
  \end{tikzpicture}
  \caption{The approximation in
    \eqref{eq:dyadic-simple-approximation-sep-17} rounds $f$ down to a
    dyadic level and truncates it at height $n$.  Below the truncation
    height, the vertical error is at most $2^{-n}$.}
  \label{fig:dyadic-approximation-sep-17}
\end{figure}

\FloatBarrier

\subsection{Nonnegative integration}

Now fix a measure space $(X,\mathcal M,\mu)$ and write
\[
  L^+
  =\{f\colon X\to[0,+\infty]:f\text{ is measurable}\}.
\]

\begin{definition}[Integral of a nonnegative simple function]
  Let $\phi\in L^+$ be simple, with standard representation
  \[
    \phi=\sum_{j=1}^{N}a_j\mathbf 1_{E_j}.
  \]
  Define
  \[
    \int_X\phi\,d\mu
      =\sum_{j=1}^{N}a_j\mu(E_j),
  \]
  using the convention $0\cdot\infty=0$.  For $A\in\mathcal M$, set
  \[
    \int_A\phi\,d\mu
      =\int_X\phi\mathbf 1_A\,d\mu.
  \]
\end{definition}

The standard representation is unique up to the ordering of its
nonempty level sets, so this definition is unambiguous.  Geometrically,
each term $a_j\mu(E_j)$ is the measure-theoretic analogue of the area of
a rectangle.

\begin{figure}[htbp]
  \centering
  \begin{tikzpicture}[x=1.35cm,y=0.85cm,>=Latex]
    \draw[->] (-0.1,0)--(7.15,0) node[right] {$x$};
    \draw[->] (0,-0.2)--(0,4.65) node[above] {$\phi(x)$};
    \fill[blue!16] (0,0) rectangle (1.4,1.2);
    \fill[green!20] (1.4,0) rectangle (3.1,3.0);
    \fill[orange!23] (3.1,0) rectangle (4.4,1.9);
    \fill[purple!17] (4.4,0) rectangle (7,3.8);
    \draw[blue!75!black,very thick]
      (0,1.2)--(1.4,1.2) (1.4,3)--(3.1,3)
      (3.1,1.9)--(4.4,1.9) (4.4,3.8)--(7,3.8);
    \foreach \x in {1.4,3.1,4.4,7} {
      \draw[gray!50] (\x,0)--(\x,4.0);
    }
    \node at (0.7,0.55) {$a_1\mu(E_1)$};
    \node at (2.25,1.5) {$a_2\mu(E_2)$};
    \node at (3.75,0.9) {$a_3\mu(E_3)$};
    \node at (5.7,1.9) {$a_4\mu(E_4)$};
    \draw[decorate,decoration={brace,mirror,amplitude=5pt}]
      (0,-0.12)--(1.4,-0.12)
      node[midway,below=6pt] {$\mu(E_1)$};
    \draw[decorate,decoration={brace,mirror,amplitude=5pt}]
      (1.4,-0.12)--(3.1,-0.12)
      node[midway,below=6pt] {$\mu(E_2)$};
  \end{tikzpicture}
  \caption{For Lebesgue measure on an interval, the simple integral is
    literally the sum of the shaded rectangular areas.  The definition
    $\sum_j a_j\mu(E_j)$ makes sense on every measure space.}
  \label{fig:simple-integral-sep-17}
\end{figure}

\FloatBarrier

\begin{definition}[Integral of a nonnegative measurable function]
  For $f\in L^+$, define
  \[
    \int_X f\,d\mu
      =\sup\left\{
        \int_X\phi\,d\mu:
        \phi\text{ is simple and }0\leq\phi\leq f
      \right\}.
  \]
  This value belongs to $[0,+\infty]$.  For $A\in\mathcal M$, define
  \[
    \int_A f\,d\mu=\int_X f\mathbf 1_A\,d\mu.
  \]
\end{definition}

Theorem~\ref*{thm:simple-approximation-sep-17} explains why the
supremum uses precisely the right class of building blocks: every
$f\in L^+$ is the pointwise increasing limit of simple functions below
it.

\begin{example}[Three measures and their integrals]
  \begin{enumerate}[label=(\roman*)]
    \item If $\mu=\delta_{x_0}$ is the Dirac measure at $x_0\in X$,
      then
      \[
        \int_X f\,d\delta_{x_0}=f(x_0)
        \qquad(f\in L^+).
      \]
      Thus a Dirac integral samples the function at one point.
    \item If $\mu=m$ is Lebesgue measure on $\bb R$, then the integral
      of a nonnegative simple function is the usual sum of rectangular
      areas, as in Figure~\ref{fig:simple-integral-sep-17}.
    \item Since $\bb Q$ is null but the fat Cantor set $K$ from
      Example~\ref{ex:fat-cantor-sep-15} has measure $1/2$,
      \[
        \int_{\bb R}\mathbf 1_{\bb Q}\,dm=0,
        \qquad
        \int_{\bb R}\mathbf 1_K\,dm=\frac12.
      \]
      Figure~\ref{fig:fat-cantor-sep-15} shows the construction of $K$.
  \end{enumerate}
\end{example}

We finish with properties that follow directly from the finite sum
defining the integral of a simple function.

\begin{proposition}[Elementary properties of the simple integral]
  \label{prop:simple-integral-properties-sep-17}
  Let $\phi,\psi\in L^+$ be simple and let $c\geq0$.
  \begin{enumerate}[label=(\alph*)]
    \item $\displaystyle\int_X c\phi\,d\mu
      =c\int_X\phi\,d\mu$.
    \item $\displaystyle\int_X(\phi+\psi)\,d\mu
      =\int_X\phi\,d\mu+\int_X\psi\,d\mu$.
    \item If $\phi\leq\psi$, then
      $\displaystyle\int_X\phi\,d\mu
      \leq\int_X\psi\,d\mu$.
    \item The set function
      \[
        \nu(A)=\int_A\phi\,d\mu,
        \qquad A\in\mathcal M,
      \]
      is a measure on $(X,\mathcal M)$.
  \end{enumerate}
\end{proposition}

\begin{proof}
  Part (a) is immediate.  For (b), write the standard representations
  as
  \[
    \phi=\sum_{j=1}^{N}a_j\mathbf 1_{E_j},
    \qquad
    \psi=\sum_{k=1}^{M}b_k\mathbf 1_{F_k}.
  \]
  The sets $E_j\cap F_k$ give a common measurable partition, and
  \[
    \phi+\psi
      =\sum_{j,k}(a_j+b_k)\mathbf 1_{E_j\cap F_k}.
  \]
  Finite additivity of $\mu$ gives
  \begin{align*}
    \int_X(\phi+\psi)\,d\mu
      &=\sum_{j,k}(a_j+b_k)\mu(E_j\cap F_k)\\
      &=\sum_j a_j\mu(E_j)+\sum_k b_k\mu(F_k)\\
      &=\int_X\phi\,d\mu+\int_X\psi\,d\mu.
  \end{align*}
  The same common partition proves (c): whenever
  $E_j\cap F_k\neq\varnothing$, the inequality $\phi\leq\psi$ implies
  $a_j\leq b_k$, so comparison term by term gives the result.

  Finally,
  \[
    \nu(A)
      =\int_X\phi\mathbf 1_A\,d\mu
      =\sum_{j=1}^{N}a_j\mu(E_j\cap A).
  \]
  Each map $A\mapsto\mu(E_j\cap A)$ is a measure.  Therefore their
  finite nonnegative linear combination $\nu$ is a measure, proving
  (d).
\end{proof}
